\section{Conclusiones}
\label{sec:conclusiones}
\begingroup
\fontsize{10.5}{13}\selectfont
\setlength{\parskip}{2pt plus 1pt minus .5pt}

La construcción determina conjuntamente \(\pi,\varphi,e\) a profundidad
arbitraria y deriva \(\alpha\) como coordenada de clausura de las regiones
y del registro dodecafásico \(K\). La misma estructura produce la
composición electrónica, su representación de espín \(1/2\) y los
proyectores de helicidad, y conserva la incidencia que conduce a Leech
y a \(V^\natural\), con el Monstruo como grupo de automorfismos.
Estos resultados proceden de las evaluaciones aditiva y multiplicativa
de APP con sus residuos y cocientes, de la orientación y los regímenes
cuadráticos de TRIT, y de la selección, el transporte y la actualización
de memoria del TPK. Sus dominios, normalizaciones y mapas de
compatibilidad determinan tanto los valores como las relaciones que
persisten entre sus realizaciones.

La exhaustividad inicial y la profundidad del refinamiento son propiedades complementarias. Las \(104\,976\) condiciones orientadas determinan \(468\) emisiones, cuya reducción contiene \(243\) palabras visibles dentro del ambiente ternario de \(729\) elementos. Sobre ese dominio se distinguen las regiones de clausura, propagación y autoescala. La región de clausura conserva cinco emisiones con multiplicidades \(432+4\cdot144=1008\). Sus realizaciones arquimedianas se determinan mediante cilindros compatibles de diámetro decreciente. Cualquier prefijo finito se obtiene mediante un cálculo finito, mientras la compatibilidad identifica el objeto completo. El retorno nonádico conserva la fase y aumenta la memoria.

La realización esturmiana precisa el orden aperiódico de ese transporte. Las separaciones primarias \(21/22\), los retornos secundarios \(6/7\) y las complejidades \(m+1\), \(9(m+1)\) y \(3(m+1)\) proceden de una misma pendiente y de sus lecturas de fase. Las relaciones de Weyl y el incremento de memoria articulan la rotación de base con el álgebra de observables y traslaciones.

La formulación mediante resolventes explicita la acumulación de memoria: la serie generatriz del registro afín se obtiene componiendo el resolvente del transporte con la serie de incrementos. La composición de tres segmentos conserva exactamente el cociclo. Cuando se elimina un subespacio auxiliar, el complemento de Schur determina la corrección del operador efectivo y, simultáneamente, la modificación de la fuente y del funcional de evaluación. Estas identidades se demuestran en el cuerpo del artículo y conservan la contribución de los estados que intervienen entre dos observaciones.

El registro dodecafásico introduce otra forma de conservación. Los canales de la trayectoria producen el vector entero firmado \(U\), y la transformación por órbitas determina
\[
 K=\tfrac14\mathcal H_{12}U,\qquad U=\mathcal H_{12}K.
\]
Estas ecuaciones expresan una equivalencia reversible entre coordenatizaciones integrales, dentro de la subred especificada. Las diferencias de pasos tres y cuatro, junto con la suma total, permiten reconstruir el mismo vector. La reversibilidad corresponde a esos registros; una lectura posterior de soportes conserva una información diferente.

La realización racional periódica
\[
 \kappa_{\mathrm{per}}=\frac{N_K}{1000^{12}-1}
\]
evalúa el registro canónico con origen y orientación fijados. Su período mínimo de doce bloques explica el denominador y distingue esta constante de la lectura finita \(N_K/1000^{12}\). Quedan separados el objeto estructural \(K\), su ecuación de reconstrucción y el número racional que lo representa. La periodicidad pertenece a esta evaluación, mientras las coordenadas regionales y sus acarreos mantienen sus propias condiciones de prolongación.

La determinación posicional de \(\alpha\) compone los registros regionales con \(K\) mediante normalización posicional. Los bloques y cocientes transportados satisfacen una igualdad local cuya suma telescópica conserva la frontera. La sección completa determina
\[
 \alpha_A=\pi+e-\varphi-4-\kappa_{\mathrm{per}}.
\]
El cuatro procede de las partes enteras regionales; la orientación determina los signos; la prolongación compatible fija el acarreo. La ecuación expresa la coordenada de compatibilidad aritmética de la clausura. Su evaluación comienza por
\[
 \alpha_A=0.007297352569283800997285105472380662999\ldots.
\]

La representación analítica correlacionada conserva el jet de orden nueve
y especifica su prolongación mediante
\(b_{10+3n}=\lambda/729^n\), \(b_{11+3n}=b_{12+3n}=0\).
El coeficiente \(\lambda\) se calcula desde las coordenadas regionales y
terminales del mismo estado, dentro de la clase geométrica declarada.
La convergencia uniforme de la función y de su derivada, la unicidad de la
raíz y los intervalos racionales compatibles con los cilindros enteros
establecen la igualdad de los prefijos de ambas representaciones a toda profundidad
(teorema~\ref{thm:alpha-dos-publicaciones-completas}). La proximidad entre
la raíz del truncamiento y \(\alpha_A\) sigue siendo una comparación finita;
la igualdad completa procede de la recurrencia y de su prueba de
compatibilidad, no de esa proximidad ni de un valor metrológico.
La coordenatización analítica es una representación del mismo punto generado,
no una segunda selección causalmente independiente.

La construcción electrónica retoma el centro de APP después de las coordenadas que utiliza. El punto fijo \((5,5)\) y las deformaciones \((8,2)\), \((2,8)\) conservan la suma y el residuo multiplicativo, mientras distinguen una unidad de cociente. Su representación matricial explica la norma \(\sqrt3/4\). La coordenada basal compone esa norma con la transferencia regional, el acoplamiento y la normalización torsional; la memoria de retorno actúa después multiplicativamente.

Dos lectores hacen explícito ese retorno: uno cuenta discrepancias cronológicas y otro compara correlaciones de ventanas dodecafásicas. Sus dominios y operaciones son diferentes, pero sobre la trayectoria electrónica ambos determinan \((80,54,6)\). La composición es
\[
 \varepsilon_e^\beta
 =\frac{\sqrt3}{4}
 \left[\exp\!\left(\frac{\varphi}{\pi^2}\right)-22\alpha^3\right]
 (1+15\Delta_4)
 R_{\mathrm{act}}^{\,80-54\alpha+6\Delta_4},
\]
cuya evaluación comienza por \(0.510998950690000808608\ldots\). La igualdad central de los triples relaciona esas representaciones sin identificar universalmente sus lectores. La escala común de acción cancela en \(R_{\mathrm{act}}\); la realización en energía o masa añade la coordenatización dimensional. El cuanto cronológico conserva su función de escala y la sección electrónica conserva su trayectoria propia.

El transporte central determina además un soporte exacto de tres diagonales
cíclicas, con \(27\) celdas. Las dos orientaciones producen registros de
eventos con el mismo histograma de valores absolutos y los mismos signos
sectoriales, pero con márgenes distintos
en la coordenatización integral reversible. La lectura armónica del registro central
concentra sus potencias en los modos \(2,6,10\), con valores \(64,16,64\).
La ponderación demostrada recupera \(\Omega=80\) y \(P_6=6\), y expresa
el exponente de retorno de la ecuación electrónica como un funcional
espectral. El histograma, la autocorrelación y el registro reversible
conservan grados de información diferentes sobre la trayectoria.

La representación espinorial del apartado~\ref{sec:el-representacion-espinorial}
explicita la información rotacional del bloque electrónico. Las dos hojas
inducen una terna hermítica y sus generadores de rotación; para cada dirección,
los proyectores espectrales separan dos líneas complejas ortogonales. La ley
de transporte distingue el cambio de signo tras una rotación completa y el
retorno espinorial tras dos. Cuando la dirección se realiza como dirección
de propagación, esta descomposición determina la helicidad. La inversión de
dirección, la conjugación de carga, la quiralidad y la holonomía nonádica
conservan sus dominios respectivos. La evaluación electrónica actúa
escalarmente sobre la representación, por lo que esta estructura rotacional
se conserva sin modificar la fórmula de masa ni sus lectores.

La prolongación excepcional desarrolla la incidencia del mismo estado. El código ternario conserva las palabras; sus soportes retienen posiciones e intersecciones. La bandera marcada de \(K\) y del registro previo al acarreo de \(\alpha\) enlaza ambas lecturas. El levantamiento binario relaciona los cuatro pétalos de Witt con una quinta octada transversal. El pegado ternario produce el retículo par unimodular, y el vecino seleccionado conduce a Leech. Con los datos algebraicos declarados, la construcción de Frenkel--Lepowsky--Meurman da \(V^\natural\), cuyo grupo de automorfismos es el Monstruo\cite{flm}. Moonshine describe después sus trazas graduadas\cite{borcherds}. La acción corresponde al álgebra obtenida; la construcción numérica y esta realización excepcional permanecen enlazadas por sus aplicaciones desde el estado común.

El contraste metrológico se realiza al final\cite{codata2018,codata2022}. \(\alpha\) coincide al redondeo con el centro de CODATA 2018 y difiere unas \(4.53\) incertidumbres estándar del centro de 2022. La expresión electrónica coincide al redondeo con el centro de 2022 y difiere unas \(4.60\) incertidumbres estándar del de 2018. La exactitud algebraica, el error de evaluación y la incertidumbre experimental tienen significados distintos. Ni la coincidencia con una fórmula ni la fecha de un ajuste deciden por sí solas su proximidad al valor físico.

El contraste histórico adicional del
apartado~\ref{sec:contraste-ventana-reciproco} utiliza la ventana finita
\(\alpha_{12}^{(0)}\) y el recíproco del centro
\(R_{2018}=137.035999084\). La diferencia en la coordenada inversa es
aproximadamente \(1.066588006664435\times10^{-32}\), con diferencia
relativa \(7.783268730800000\times10^{-35}\); en la coordenada directa
es aproximadamente \(-5.679725607012965\times10^{-37}\).
Esta proximidad corresponde al truncamiento decimal del recíproco
impreso. Su precisión aritmética, la comparación con el centro tabulado
de \(\alpha\) y la incertidumbre experimental conservan sus dominios
respectivos.


La dirección \(u_K=P_3P_{11}K\) amplía esta descripción de la función
de \(K\). Su norma exacta positiva demuestra que el registro selecciona
una recta, y el proyector ortogonal complementario realiza la reducción
\(11\to10\). La reducción elimina una dirección especificada por el
registro; el mapa coordinado conserva además la pantalla reticular
\(26\to24\) y entrelaza el intercambio de radios recíprocos. Cada
factor mantiene su propio dominio y su información conservada.

La construcción excepcional admite asimismo dos realizaciones por
orbifolds, de órdenes \(2\) y \(3\), sobre el mismo vecino reticular.
Las identificaciones con \(V^\natural\) proporcionan un isomorfismo
que conserva el vacío, el vector conforme, el producto de vértice y
la gradación. Las trazas de Fricke y el automorfismo dual del orbifold
describen operaciones concretas sobre estas realizaciones. El exponente
\(-1/12\) del cociente de funciones eta de nivel tres tiene aquí
una residencia precisa, distinta de una suma ordinaria divergente.
Estas aplicaciones enlazan las dos realizaciones mediante morfismos
que conservan sus datos algebraicos y sus identificaciones clásicas.

% BEGIN S27_FRAMING
El cociclo de acarreo determina cómo se conserva la evaluación entera
al componer columnas regionales. Las identidades
\(p+e+f=q+3c\) y \(p+e-f=a+3h\) recuperan las dos combinaciones
enteras desde sus residuos y cocientes. La tabla completa muestra
que \((0,0,0)\) y \((1,2,0)\) comparten \(q=a=0\), mientras sus
cocientes son, respectivamente, \((c,h)=(0,0)\) y \((1,1)\).
Las representaciones polinómicas y los rangos \(4,5,5\) precisan
la contribución algebraica de esos acarreos. Las nueve fibras de tres
ternas del mapa residual \((p,e,f)\mapsto(q,a)\) conservan explícita
la libertad de reparto entre las entradas, y las recurrencias
transportan las evaluaciones conjuntas y firmadas en el dominio
entero declarado.

La lectura incidencial distingue las tres estrellas regionales por
las firmas \((3,1,0)\), \((1,0,3)\) y \((1,1,2)\), con sus doce
completados explícitos. El censo global distribuye las \(495\)
cuaternas en seis clases, y la restricción regional clasifica las
\(54\) hexadas pertinentes. La sucesión de filtros sobre el grupo
construido prueba que el marco ordenado
\((H^-_\alpha,H_e,H_\varphi,H_{\rm APP})\) tiene estabilizador
trivial. Esta rigidez determina la libertad residual de su acción
sobre las doce posiciones. La recuperación numérica de \(K\)
conserva sus propias ecuaciones reversibles y el registro entero
completo. Quedan así articuladas dos funciones de la misma
genealogía: magnitudes, residuos y cocientes en la evaluación
aritmética; soportes, polarización y estabilizadores en la
realización incidencial.

% END S27_FRAMING

% BEGIN R27_FRAMING_CONCLUSION_ES
Las pruebas del emisor y de su prolongación precisan los grados de
información que conserva esta genealogía. Las dos firmas determinan los
\(42\) bloques decimales posibles de una ventana; seis emisiones testigo
establecen el ciclo de empalmes regionales. Un histograma común puede
acompañar a correlaciones temporales distintas, de modo que el archivo
ordenado conserva una distinción operativa adicional. Los controles de
calibre y de frontera identifican, respectivamente, la normalización del
observable y la dependencia respecto del sentido de acarreo, del extremo
terminal y de la posición de los coeficientes.

En APP, las leyes del defecto conservan explícitos módulo, soporte y
dimensión. La restricción del operador agregado a una subretícula de
índice dos produce \((-1)\oplus3H\), con \(\det H=1\). La recurrencia
de Pell, con los datos iniciales fijados, determina las potencias de
\(H\), que conservan norma uno. La codificación nonádica conserva
los seis trits mediante tres pares y transporta la suma con acarreo.
La correlación cuadrática da \(C^2=5I_6\); la saturación, la neutralidad
dual y la orientación fijan por etapas la frontera de prolongación.
Finalmente, la identidad de nivel tres demuestra
\(J-(t_3+12)\in3^9q\mathbb Z[[q]]\), con exponente uniforme óptimo.
Estos resultados enlazan las operaciones finitas, sus leyes generales
demostradas y la realización excepcional, manteniendo el dominio propio
de cada afirmación.
% END R27_FRAMING_CONCLUSION_ES

% BEGIN R28_FRAMING_CONCLUSION_ES
Los cinco cierres regionales alcanzan el mínimo ocho en el problema
finito de cobertura de las cuatro aristas prescritas. El grafo conserva
las etiquetas de las emisiones y el cociente conserva sus incidencias
entre clases de fase. La inversión orientada hace impar al lector
\(\nu_{120}\) y mantiene par a \(\nu_{270}\); estas identidades
explicitan la respuesta de la memoria al cambio de orientación.
% END R28_FRAMING_CONCLUSION_ES
\par % S27_LAYOUT_CONCLUSION; REV03: continuous conclusion pagination.
La completación cúbica y el transporte aperiódico precisan dos sentidos
de conservación que se articulan en el desarrollo. La primera conserva
la unidad bajo proyección de medidas normalizadas y distingue el cubo
inicial, la completación perforada de \(999\) elementos y la restitución
del ápice. El segundo conserva el registro de trayectoria durante el
retorno de fase. La monodromía y el modelo matemático de cristal temporal
aperiódico expresan este segundo mecanismo mediante aplicaciones y
observables, no sólo mediante una imagen helicoidal. La exponencial
trítica, por su parte, unifica las realizaciones elíptica, parabólica e
hiperbólica sin suprimir sus relaciones cuadráticas respectivas.

La constante de Planck reducida, \(\hbar\), ocupa una posición explícita
en esta arquitectura. El cociente determinantal de orden \(16\) actúa
en la norma que determina la década \(-34\); el carácter \(H_5\) y su
continuación regional determinan las dos secciones de acción cuyo
cociente entra en la coordenada electrónica. La comparación de la
sección anterior con \(h/(2\pi)\) muestra una diferencia relativa del
orden de \(10^{-15}\), y la sección retornada conserva la compensación
que utiliza el electrón. La escala de acción enlaza así la norma
determinantal y el carácter regional con el cociente electrónico: la
compensación pertenece a la construcción exacta, mientras que la proximidad
de la sección anterior a \(h/(2\pi)\) es una comparación numérica y no
una igualdad exacta.


La consecuencia conceptual se concentra en la articulación de \(K\),
\(\alpha\) y la sección electrónica. La reconstrucción reversible de
\(K\) conserva datos que intervienen tanto en la normalización posicional
como en la incidencia de frontera. La composición electrónica utiliza las
coordenadas previamente construidas y mantiene su centro, sus transportes
y su realización dimensional. La prolongación excepcional conserva otras
relaciones del mismo estado bajo representaciones sucesivas. Las conexiones
entre estos objetos forman, por tanto, parte del resultado junto con los
valores de sus coordenadas.

Esta organización determina también el orden de la contrastación. Cuando
una derivación fija un valor sin recibirlo como argumento, la medición
somete a prueba la consecuencia obtenida. La correspondencia con un
observable físico exige, además, identificar la magnitud, la normalización
y la realización dimensional empleadas. En ese sentido preciso, la tesis
examina el tránsito de una descripción parametrizada a la explicación de
las relaciones que determinan sus parámetros. El alcance de cada resultado
es el de las demostraciones expuestas: estructura, ecuación y valor
constituyen aspectos inseparables de la determinación que se contrasta.

La unicidad del registro conserva dos censos prospectivos y sus
operaciones propias (apartado~\ref{act21:sec:relacion-censos}).
La cadena \(40320\to21902\to19446\to79\to1\) culmina en el
superviviente heterotípico de la cara excepcional; la cadena
\(196\cdot197\cdot196\to331\to2\to1\) culmina en la selección
orientada desde catálogos regionales. Las poblaciones intermedias
conservan el significado de cada filtro. La composición de las salidas
seleccionadas con la inversión integral articula determinación
prospectiva y recuperación reversible.


Al cerrar su conferencia Nobel de 1946, Pauli vinculó la explicación de la electricidad con una teoría capaz de \emph{“determine the value of the fine-structure constant”}\cite{pauli1946}: determinar el valor de la constante de estructura fina. La respuesta matemática obtenida aquí es la identidad de clausura de \(\alpha\), su representación analítica compatible y su composición con la estructura electrónica. La generación regional, la reconstrucción reversible de \(K\) y los transportes de incidencia hacen de los valores y de sus relaciones consecuencias de una misma construcción. La metrología contrasta después sus realizaciones físicas; no fija retrospectivamente los estados ni los coeficientes que las determinan.



\par\endgroup
